\documentclass[12pt]{article}

% Layout.
\usepackage[top=1in, bottom=0.75in, left=0.75in, right=0.75in, headheight=1in, headsep=6pt]{geometry}

% Fonts.
\usepackage{mathptmx}

\usepackage[scaled=0.86]{helvet}
\renewcommand{\emph}[1]{\textsf{\textbf{#1}}}

% Misc packages.
\usepackage{amsmath,amssymb,latexsym,amsthm}
\usepackage{graphicx,hyperref}
\usepackage{array}
\usepackage{xcolor}
\usepackage{multicol}
\usepackage{tabularx,colortbl}
\usepackage{enumitem}
\usepackage{soul,multicol}
\usepackage{setspace}


\hypersetup{
    colorlinks=true,
    linkcolor=blue,
    filecolor=magenta,      
    urlcolor=blue,
    pdftitle={Proofs Worksheet},
    pdfpagemode=FullScreen,
    }

\def\mailto#1{\href{mailto:#1}{#1}}

%% special math symbols
\newcommand{\N}{\mathbb{N}}
\newcommand{\Z}{\mathbb{Z}}
\newcommand{\R}{\mathbb{R}}
\newcommand{\Q}{\mathbb{Q}}
\newcommand{\sse}{\subseteq}
\newcommand{\mcp}{\mathcal{P}}
\newcommand{\ol}[1]{\overline{#1}}

%other specials
\newcommand{\be}{\begin{enumerate}}
\newcommand{\ee}{\end{enumerate}}
\newcommand{\se}{\subseteq}
\newcommand{\spe}{\supseteq}
\newcommand{\ds}{\displaystyle}
\newcommand{\lcm}{\text{lcm}}
%\newcommand{\gcd}{\text{gcd}}

% Paragraph spacing
\parindent 0pt
\parskip 6pt plus 1pt
\def\tableindent{\hskip 0.5 in}
\def\ts{\hskip 1.5 em}

%header
\usepackage{fancyhdr}
\pagestyle{fancy} 
\lhead{\large\sf\textbf{MATH 265: Introduction to Mathematical Proofs}}
\rhead{\large\sf\textbf{Worksheet 20: Ch 11}}

\newcommand{\localhead}[1]{\par\smallskip\textbf{#1}\nobreak\\}%
\def\heading#1{\localhead{\large\emph{#1}}}
\def\subheading#1{\localhead{\emph{#1}}}

\newenvironment{clist}%
{\bgroup\parskip 0pt\begin{list}{$\bullet$}{\partopsep 4pt\topsep 0pt\itemsep -2pt}}%
{\end{list}\egroup}%

\begin{document}
\begin{enumerate}
\item Quick Review
\vspace*{-0.5cm}
	\begin{enumerate}
	\item $R$ is a relation on $A$ if
	
	\item Suppose $R$ is a relation on the set $A.$
		\begin{enumerate}
		\item We say $R$ is \emph{symmetric} if \\
		\item We say $R$ is \emph{reflexive} if \\
		\item We say $R$ is \emph{transitive} if \\
		\end{enumerate}
	\item What properties do the following relations have?
		\begin{enumerate}
		\item Let $R$ be a relation on $\N \times \N$ such that $(a,b) R (c,d)$ if $ad=bc.$
		\vfill
		\item Let $R$ be a relation on $A=\Z$ defined by $a \;R\; b$ if $a \equiv b \pmod 3.$ 
	\vfill
		\item Suppose $R$ is a relation on $\mathcal{P}(A)$ where $A=\{0,1,2,3,4\}$ defined as $X \: R \: Y$ if $X \cap Y \not = \emptyset.$\\
		\vfill
		\end{enumerate}
	\end{enumerate}
\item A relation $R$ on set $A$ is an \textbf{equivalence relation} if \\
	\vfill
	\vfill
\item Build your own equivalence relation on $A=\{a,b,1,5\}.$
\vfill
\vfill
\newpage
\item For each relation below, quickly confirm the relation \emph{is} an equivalence relation, then identify its equivalence classes using the square bracket notation \emph{in two different ways} and by describing them as sets. Can you describe the partition of $A$ that is produced?
	\begin{enumerate}
	\item Let $R$ be a relation on $A=\Z$ defined by $a \;R\; b$ if $a \equiv b \pmod 6.$
	\vfill
	\item Let $R$ be a relation on the set of all polynomials with real coefficients defined by $p(x) R q(x)$ if the degree of $p(x)$ equals the degree of $q(x).$
	\vfill
	\item Let $R$ be the set of differentiable functions $f:\R \to \R$ defined by  by $f(x) R g(x)$ if $f'(x)=g'(x).$
	\vfill
	\end{enumerate}
\item Suppose $R$ and $S$ are two equivalence relations on $A$. Answer the following questions \emph{rigorously}. 
	\begin{enumerate}
	\item Is the set $R \cup S$ an equivalence relation on $A$?
	\vfill
	\item Is the set $R \cap S$ an equivalence relation on $A$?
	\vfill
	\vfill
	\end{enumerate}	
\end{enumerate}
\end{document}










